\begin{lemma}[Perfect-or-bad shift dichotomy]
For every relation $r$ and every locally constant multi-sequence $h$, there
is an infinite restriction on which $h$ is either perfect or bad.
\end{lemma}
\begin{proof}
Define the Boolean multi-sequence
\[
c_h(x)=\begin{cases}1&\text{if }r(h(x),h(\mathsf{ShiftMap}(x))),\\
0&\text{otherwise.}\end{cases}
\]
The shift changes only the tail of an increasing enumeration, so a finite
prefix deciding $h$ at $x$ together with a finite prefix deciding $h$ at
$\mathsf{ShiftMap}(x)$ gives a finite prefix deciding $c_h$ at $x$.
Thus $c_h$ is locally constant.  Apply \noderef{DecidingFront} and
\noderef{DecidingValue} to $c_h$: they give a deciding front $F^{c_h}$ and
a value map $v:F^{c_h}\to\{0,1\}$ whose value on the front member chosen
by an enumeration is $c_h$ at that enumeration.

Let $S=\{s\in F^{c_h}\mid v(s)=1\}$.  By \noderef{NashWilliams}, there
are a sub-front $F'\subseteq F^{c_h}$ on an infinite base $Y$ such that
either $F'\subseteq S$ or $F'\cap S=\emptyset$.  Choose the increasing
enumeration $z$ of $Y$ using \noderef{InfiniteSetEnumeration}.  For each
$x\in\mathsf{IncSeq}$, the front member selected by the range of
$z\circ x$ belongs to $F'$, so the value equation for $v$ makes
$c_h(z\circ x)$ constantly $1$ in the first case and constantly $0$ in
the second.

In the first case, the definition of $c_h$ gives

\[
r(h(z\circ x),h(\mathsf{ShiftMap}(z\circ x))).
\]

By \noderef{RestrictionShift}, the second argument is
$h(z\circ\mathsf{ShiftMap}(x))$.  By
\noderef{MultiSequenceRestrict}, these are respectively the values of
$\mathsf{MultiSequenceRestrict}(h,z)$ at $x$ and at
$\mathsf{ShiftMap}(x)$; hence \noderef{PerfectMultiSequence} holds.  In the
second case the same calculation gives the negation of this relation for
every $x$, which is \noderef{BadMultiSequence}.  This proves the displayed
perfect-or-bad alternative.
\end{proof}
